% Transcription — fonds Alexandre Grothendieck, shelfmark 5, batch 2 (pages 21-40).
% Established from the facsimile archives/batches/5/batch-02.pdf (a mirror of
% https://grothendieck.umontpellier.fr/5.pdf), per the skill
% transcribe-grothendieck, read off the tiles archives/tiles/5/p21 … p40. The
% batch file carries no cover sheet and is cut from the folder starting at page
% 21: PDF page 1 is archivists' page 21, PDF page N is archivists' page N+20.
% Checked against the pencilled numbers bottom-left — "31", "32", "33", "34",
% "35", "36", "37", "38", "39", "40" read clearly on the corresponding leaves at
% tile resolution; the alignment holds end to end.
%
% Pass: Fable 5.1 (claude-fable-5-1), 2026-09-19 — first pass, unchecked against the pages by a human.
%
% What the batch is. Sixteen leaves in his hand and four typed versos. Pages 21
% to 26 are a numbered run, § 1 to § 3 in his numbering, on pseudo-rings with
% divided powers: the axioms 1.1–1.6, the divided-power algebra Γ(M) of a
% module, the universal divided-power envelope Γ(A, I) of an ideal, its
% quotients, nilpotent divided powers and the log/exp isomorphism. Pages 27 to
% 30 are § 4, the De Rham site (X/S)_DR, its stratified variant, the square of
% topoi Zar / DR_st / DR / Strat / Cris, and the Leray spectral sequence that
% computes H*(X_DR st, O) from the infinitesimal neighbourhoods of the diagonal.
% Pages 31, 32 and 34 return to the envelope: generators Γ^n(x) and relations,
% the p-adic simplification (only the Γ^{p^ν} are needed when the primes ≠ p
% are invertible), the coefficient c_n, and the case of a principal ideal (λ).
% Page 36 is a short list of k-algebras for which the Poincaré lemma holds, with
% a reference. Pages 38 and 40 are the perfect-residue-field case: the
% operations π_m are determined by π_p alone, the coefficients c_n and d_m
% have p-adic valuation computed, and a construction of π_n from the p-adic
% digits of n is begun. Pages 33, 35, 37 and 39 carry no mathematics of his
% and are absent.
%
% The run does not start on page 21: § 1 opens there, so batch 1 presumably
% closes something else. Page 32 breaks off mid-sentence and page 34 completes
% it — page 33 is a typed verso between them. Page 40 ends with a question and
% a construction begun; batch 3 opens on whatever follows.
%
% Notation fixed for the batch. His divided-power operation is written three
% ways on these leaves and each is kept: π^n(x) = x^{(n)} in § 1–3 (pages
% 21–26), Γ^n(x) — capital gamma as an operation, not the functor — on pages
% 31, 32 and 34, where it is a generator of a polynomial ring, and π_m (index
% below) on pages 38 and 40. The one-off γ^n(x) on page 31 is kept as γ. The
% functor Γ^n(M), Γ^+(M), Γ_Λ(M) is set as \Gamma too; the context separates
% them. Sites and topoi: DR, DR_st, Strat, Cris in \mathrm, the topos with a
% tilde as he draws it, \widetilde{\mathrm{DR}}(X). The diagonal neighbourhoods
% are Δ^{(ν)}_{X/S}(i) and Δ^{(ν)}_{X/S}(i, j) as written. His « ps-anneau »
% and « puiss. div. » are kept as abbreviations. Number sets in \mathbb.
%
% The argument, page by page:
%
%   21    pseudo-ring; divided powers π^n on a pseudo-ring J, axioms 1.1–1.5;
%         when J is an ideal of A, 1.2 with π^0 = 1; 1.5 iterated (1.5 bis)
%         and n! π^n(x) = x^n (1.5 ter).
%   22    A-algebra with divided powers, axiom 1.6 (automatic over Z);
%         divided powers on an ideal vs on the underlying algebra. § 2: Γ*(M)
%         and the extension of x ↦ x^{(n)} to π^n : Γ^+(M) → Γ^+(M), reducing
%         to a homogeneous polynomial map and the formula for (x^{(p)})^{(n)}.
%   23    Γ^+(M) is the universal divided-power A-algebra under M. § 3: the
%         universal (B, J, u) with u(I) ⊂ J for an ideal I ⊂ A, as a quotient
%         of Γ^+_A(I) by relations.
%   24    the quotient J' by the (x.y)^{(n)} − x^n * y^{(n)}; the extension of
%         A by J' along J → J'; N.B. on Λ-algebra structure.
%   25    special case A = Sym(M), J = Sym^+(M), giving Γ_Λ(M). Killing a
%         family (x_i) ⊂ J: the smallest ideal stable under the π^n.
%   26    quotients (A/K, J/K) keep divided powers; nilpotent divided powers,
%         nil-ideals; log and exp inverse to each other between 1+J and J.
%   27    § 4: objects (U, U', i, π) of the De Rham site; Zariski topology;
%         the stratified site (a retraction exists locally); the topoi.
%   28    the square Zar → DR_st → Strat over DR → Cris; vertical arrows are
%         equivalences for X formally smooth, horizontal ones in char. 0;
%         fibre products exist in DR and DR_st, doubtful in Cris and Strat;
%         the Leray spectral sequence for X̃ ∈ DR_st.
%   29    X̃^{ν+1} is ind-representable by the neighbourhoods Δ^{(ν)}(i, j);
%         H^q(DR_st/(U, U'), O) ≃ H^q(U', O_U'), « à vérifier ».
%   30    hence for X affine H^q(X̃^{ν+1}, O) = 0 for q > 0 and H^n(X_DR st, O)
%         is the cohomology of ν ↦ H^0(X̃^{ν+1}); in general H*(X_Zar, E*).
%   31    B = A[Γ^n(x)] / relations, with the four relations; J generated by
%         the Γ^n(x); a struck block; if primes ≠ p are invertible, v_p(n!)
%         from the p-adic digits and γ^n(x) = c_n ∏ γ^{p^i}(x)^{a_i}.
%   32    c_n as the prime-to-p part of ∏(p^i!)^{a_i}/n!; B generated by the
%         [x, ν] = Γ^{p^ν}(x); the three relations among them; marginal
%         Γ^{p^ν}(x)^p = (p^{ν+1})!/((p^ν!)^p p!) Γ^{p^{ν+1}}(x); I = (λ).
%   34    the relations for I = (λ): Γ^m Γ^n − binomial Γ^{m+n} and
%         (x^n − y^n)Γ^n(λ) when (x − y)λ = 0; N.B. reducing to (**); the
%         p-adic version and a boxed, overwritten formula for Γ^{p^ν}(λ)^e.
%   36    Poincaré lemma for a list of k-algebras, char. 0; a remark on
%         dx = 0 ⇒ x ∈ m^3 for local A with trivial residue extension;
%         a reference to Ann. Éc. Norm. Sup. 65 or 66 and J. de Liouville.
%   38    perfect residue fields, p > 0: a) π_m π_n = C(m+n, n) π_{m+n} and
%         its inversion when m ≡ 0 (p), 0 ≤ n ≤ p−1; b) π_m ∘ π_n and the
%         valuations of c_n = (pn)!/((n!)^p p!) and d_m = (pm)!/((p!)^m m!);
%         π_n is known once π_p is.
%   40    the two conditions on π = π_p; whether the π^{(n)} can be recovered
%         from π; a struck inductive construction; π_n(x) = c_n ∏ π^i(x)^{a_i}
%         with c_n ∈ Q ∩ Z_p^*, and the question whether c_n ≡ 1 mod p.
%
% Reading notes. The hand is a fast draft throughout, with many interlinear
% substitutions and several blocks struck through; formulas carry, connective
% prose often does not. Two slips are on the page and flagged, not repaired:
% the log series on page 26 (bounds and sign) and the valuation identity on
% page 31 (v_p(n) for v_p(n!)). The boxed formula on page 34 is overwritten
% past reconstruction and is transcribed as it reads, with a note. Page 36's
% « m^3 » is a clean reading. The typed versos (33: a typescript in French on
% quasi-functions on an abelian variety, paginated 41; 35: a typescript
% paginated 39 bis; 37: a typescript page 18 of an EGA-style text; 39: an
% upside-down course programme) are skipped as unrelated to the notes.

\documentclass[11pt,a4paper]{article}
% Le préambule commun à toutes les transcriptions du fonds Grothendieck.
%
% Il tient en une idée : **l'incertitude se compose, elle ne se résout pas.**
% Une transcription qui lisse un mot douteux a détruit la seule chose pour
% laquelle on la faisait — un lecteur doit pouvoir savoir, à chaque endroit,
% ce qui était sur le feuillet et ce qui a été ajouté. Les six macros qui
% suivent sont donc l'appareil critique, et non de la décoration.
%
% Les mêmes commandes sont reconnues par `scripts/render.mjs`, qui produit la
% vue de lecture du site. Une macro ajoutée ici doit l'être aussi là-bas,
% faute de quoi le rendu échouera bruyamment — ce qui est voulu : un rendu
% silencieusement faux est pire qu'un rendu absent.

\ProvidesPackage{grothendieck}[2026/08/08 Transcriptions du fonds Grothendieck]

\usepackage{amsmath,amssymb,amsthm}
\usepackage{mathrsfs}
\usepackage{stmaryrd}
% Les diagrammes commutatifs sont partout dans ce fonds : c'est la langue
% même de la géométrie algébrique de Grothendieck, et une transcription qui
% les rendrait en prose ne transcrirait rien.
\usepackage{tikz-cd}
% Français par défaut — c'est la langue des feuillets — mais l'anglais est
% chargé aussi : la traduction et le résumé sont des documents anglais, et la
% typographie française (espaces avant les deux-points) y serait une faute.
% Ces documents basculent par \selectlanguage{english} en tête de corps.
\usepackage[english,french]{babel}
\usepackage{fontspec}
\usepackage{geometry}
\geometry{a4paper,margin=25mm,marginparwidth=28mm}
\usepackage{marginnote}
\usepackage{xcolor}
% ulem, not soul. `soul`'s \st and \ul are built on a character-by-character
% scan that XeLaTeX's font machinery breaks: the document fails with an
% undefined control sequence rather than degrading. ulem does the same two jobs
% and is Unicode-engine safe. `normalem` stops it hijacking \emph, which the
% transcriptions use for Grothendieck's own emphasis.
\usepackage[normalem]{ulem}
\usepackage{hyperref}
\hypersetup{colorlinks=true,linkcolor=black,urlcolor=black,pdfborder={0 0 0}}

% --- Métadonnées du lot -----------------------------------------------------
% Reprises telles quelles par le script de rendu, qui les affiche en tête de
% la vue de lecture. Elles fixent ce que le fichier prétend couvrir : sans
% elles, une transcription flotte sans point d'attache dans un fonds de seize
% mille feuillets.

\newcommand{\folder}[1]{\gdef\grothFolder{#1}}
\newcommand{\batch}[1]{\gdef\grothBatch{#1}}
\newcommand{\leaves}[2]{\gdef\grothFirst{#1}\gdef\grothLast{#2}}
\newcommand{\foldertitle}[1]{\gdef\grothFoldertitle{#1}}
\gdef\grothFolder{?}\gdef\grothBatch{?}\gdef\grothFirst{?}\gdef\grothLast{?}\gdef\grothFoldertitle{}

% --- Le résumé --------------------------------------------------------------
%
% Il ouvre la lecture modernisée, sous le titre « Résumé », et il est composé
% en petit. Ce n'est pas un ornement : c'est le seul passage du document qui
% ne soit pas des mathématiques, et un lecteur doit pouvoir le reconnaître
% avant de l'avoir lu — puis le sauter s'il connaît déjà le sujet, ou s'y
% arrêter s'il ne le connaît pas. Le filet qui le suit marque la reprise du
% corps.

\newenvironment{resume}
  {\small\setlength{\parskip}{0.45em}}
  {\par\medskip\hrule\medskip}

% --- Mots-clés ---------------------------------------------------------------
%
% En anglais, en clôture du résumé de la lecture modernisée : le vocabulaire
% moderne sous lequel chercher ce que les feuillets construisent. Le manifeste
% du site (`npm run manifest`) les extrait pour étiqueter la cote entière —
% cette ligne est la source unique des tags, il n'y a pas de fichier à côté.
\newcommand{\keywords}[1]{\par\smallskip\noindent
  {\footnotesize\textsc{Keywords} --- \textit{#1}}\par}

% --- L'appareil critique ----------------------------------------------------

% Le numéro de feuillet, en marge. C'est l'ancre qui permet de régler un
% désaccord sur une formule en regardant *un* feuillet plutôt qu'une chemise
% de six cents. Le site s'en sert aussi pour tourner les pages du fac-similé
% quand on fait défiler la transcription.
\newcommand{\leaf}[1]{%
  \marginnote{\footnotesize\color{gray!70}\textsf{#1}}%
  \label{leaf:#1}\ignorespaces}

% Illisible : on ne devine pas. Un mot inventé qui se lit comme les autres est
% le pire résultat possible.
\newcommand{\ill}{\textcolor{red!60!black}{[\,\ldots\,]}}

% Lecture douteuse : le mot est donné, et signalé comme incertain.
\newcommand{\uncertain}[1]{\uline{#1}}

% Ajout de l'éditeur — une lettre manquante, un mot sous-entendu.
\newcommand{\add}[1]{\textcolor{blue!50!black}{[#1]}}

% Biffé par Grothendieck lui-même. Ce qu'il a rayé fait partie du document :
% une rature dit souvent où la pensée a changé de direction.
\newcommand{\struck}[1]{\sout{#1}}

% Note du transcripteur, sur l'état matériel du feuillet.
\newcommand{\note}[1]{\par\smallskip{\small\color{gray!60!black}\textit{[#1]}}\par\smallskip}

% Note marginale de Grothendieck, distincte du corps du texte.
\newcommand{\marginal}[1]{\par\smallskip\noindent\hspace{1em}%
  {\small\color{gray!60!black}#1}\par\smallskip}

% --- Titre ------------------------------------------------------------------

\AtBeginDocument{%
  \begin{center}
    {\large\bfseries Fonds Alexandre Grothendieck --- cote n\textsuperscript{o}~\grothFolder}\\[2pt]
    {\itshape\grothFoldertitle}\\[4pt]
    {\small feuillets \grothFirst--\grothLast\ (lot \grothBatch)}\\[2pt]
    {\footnotesize\color{gray!60!black}Transcription établie d'après le fac-similé numérisé
     par l'Université de Montpellier.\\
     Les lectures incertaines sont soulignées, les illisibles notées [\,\ldots\,],
     les ajouts entre crochets.}
  \end{center}
  \vspace{1em}}

\endinput


\folder{5}
\batch{2}
\pages{21}{40}
\dating{1965-1969}
\watermark{Édition de démonstration}
\foldertitle{Cristaux et cohomologie de De Rham : généralités, correspondances Deligne et Berthelot (1965-1969) : notes manuscrites (s.d.), lettre (1965).}

\begin{document}

\note{Le § 1 s'ouvre en tête du feuillet 21 ; ce qui le précède dans le
dossier n'est pas dans ce lot.}

\section{Pseudo-anneaux à puissances divisées}

\page{21}
\textbf{§ 1.} Pseudo-anneau $=$ groupe abélien à loi de composition interne
associative bilinéaire. Nous sous-entendrons par la suite la commutativité.

Pseudo-anneau $J$ \emph{à puissances divisées}. C'est un ps-anneau muni
d'applications \struck{satisfaisant}
\[
  \pi^{n} : x \mapsto x^{(n)} = \pi^{n}(x) \qquad n \in \mathbb{N}^{+}
\]
satisfaisant les conditions
\begin{enumerate}
  \item[1.1.] $\pi^{1} = \mathrm{id}$ \quad i.e. $x^{(1)} = x$
  \item[1.2.] $\pi^{n}(x+y) = \struck{\ill{}}\; \pi^{n}(x) + \sum_{\substack{p,q \geqslant 1 \\ p+q = n}} \pi^{p}(x)\,\pi^{q}(y) + \pi^{n}(y)$
  \item[1.3.] $\pi^{n}(xy) = x^{n}\,\pi^{n}(y)$
  \item[1.4.] $\pi^{m}(\pi^{n}(x)) = \dfrac{(mn)!}{(n!)^{m}\, m!}\, \pi^{mn}(x)$
  \item[1.5.] $\pi^{m}(x)\,\pi^{n}(x) = \dfrac{(m+n)!}{m!\, n!}\, \pi^{m+n}(x)$
\end{enumerate}
\note{en 1.2, un pâté d'encre après le signe $=$ recouvre un premier terme biffé}

Supposons que $J$ soit un idéal d'un anneau $A$, avec la structure de
pseudo-anneau induite. Alors la relation 2. peut s'écrire
\[
  \pi^{n}(x+y) = \sum_{p+q=n} \pi^{p}(x)\,\pi^{q}(y) \qquad x, y \in J
\]
si on pose $\pi^{0}x = 1$, $\pi^{p}(x) = 0$ si $p < 0$. Formule 5.
\uncertain{itérée} :
\[
  \text{(1.5.bis)}\qquad \pi^{p_{1}}(x)\,\pi^{p_{2}}(x) \cdots \pi^{p_{r}}(x)
  = \frac{(p_{1} + \cdots + p_{r})!}{p_{1}! \cdots p_{r}!}\; \pi^{p_{1} + \cdots + p_{r}}(x)
\]
en particulier
\[
  \text{1.5.ter}\qquad n!\; \pi^{n}(x) = \struck{\ill{}}\; x^{n}
\]
\note{le numéro « 1.5.ter » est écrit par-dessus un premier numéro, et un exposant biffé suit $\pi^{n}(x)$}

\struck{\ill{}} $A$ est un anneau (commutatif), une $\uncertain{\Lambda}$-algèbre

\page{22}
à puissances divisées est un \supplied{pseudo-}anneau à puissances divisées, muni
d'une structure de $A$-algèbre, satisfaisant
\note{« pseudo- » est ajouté en interligne au-dessus de « anneau »}
\[
  \text{1.6.}\qquad \pi^{n}(\lambda x) = \lambda^{n}\, \pi^{n}(x) \qquad \lambda \in A,\ x \in \struck{A}.
\]
\marginal{N.B. C'est \uncertain{tjs} vérifié si $A = \mathbb{Z}$}

Si $J$ est un idéal de $A$, c'est une $A$-algèbre, \struck{\ill{}} \uncertain{quelle} structure :
puissances divisées sur l'idéal, \struck{\ill{}} $J$ une structure de puissances
divisées sur \struck{\ill{}} \add{la} $A$-algèbre sous-jacente, i.e. une structure de
p. div. sur le pseudo-anneau sous-jacent, satisfaisant 6.
\note{deux lignes biffées ; « quelle structure : puissances divisées sur » et « $A$-algèbre » sont en interligne}

\textbf{§ 2.} Soit $A$ un anneau, et $M$ un $A$-module. Considérons
\add{la $A$-algèbre}
\[
  \Gamma^{*}(M) = \coprod_{n \geqslant 0} \Gamma^{n}(M),
\]
\ill{} les applications $\pi^{n}_{(1)} : x \mapsto x^{(n)}$ de $M$ dans
$\Gamma^{n}M$. On va les prolonger en des applications
\[
  \pi^{n} : \Gamma^{+}(M) \longrightarrow \Gamma^{+}(M) \qquad x \mapsto x^{(n)}
\]
induisant
\[
  \pi^{n}_{(p)} : \Gamma^{p}(M) \longrightarrow \Gamma^{np}(M) \qquad \struck{x \mapsto x^{(n)}}
\]
en notant que ceci revient à donner une « application polynomiale homogène »
de degré $\uncertain{np}$
\[
  x \mapsto \pi^{n}_{(p)}(x^{(p)}) : \underline{M} \longrightarrow \underline{\Gamma^{np}(M)}
\]
\note{les deux termes de la flèche sont soulignés d'un trait ondulé ; sous $\pi^{n}_{(p)}(x^{(p)})$ il écrit « $= (x^{(p)})^{(n)}$ »}
et \ill{} \uncertain{pose}
\[
  (x^{(p)})^{(n)} = x^{(pn)}\, \frac{(pn)!}{(p!)^{n}\, n!}
\]
\page{23}
On vérifie que les $\pi^{n} : \Gamma^{+}(M) \to \Gamma^{+}M$ ainsi
construits satisfont les conditions 1.1. à 1.6. ($\Gamma^{+}(M)$ considéré
comme un idéal de $\Gamma(M)$, ou comme une $A$-algèbre).

$\Gamma^{+}(M)$ devient ainsi une $A$-algèbre à puissances divisées. Elle est
universelle \uncertain{muni} d'un hom de $A$-modules
\[
  M \longrightarrow \Gamma^{+}(M) ;
\]
ce qui donne pour un tel \ill{} un rôle \emph{universel}.

\textbf{§ 3.} Soit $A$ un anneau, et $I$ un idéal de $A$. On s'intéresse
à un anneau $B$, muni d'un idéal $J$ \emph{à puissances
divisées}, et $x$ un hom d'anneaux $u : A \to B$ tel que
$u(I) \subset J$. On demande un tel \uncertain{système} qui
soit le \uncertain{fourni}, \emph{universel}. On peut
l'\struck{\ill{}} $A$-algèbre \note{« $A$-algèbre » en interligne au-dessus d'un mot biffé}
$\Gamma^{*}(A, I)$ engendré par des éléments $\pi^{n}x$ \struck{$n \in \mathbb{N}^{*}$}
$n \geqslant 2$, $x \in I$ \note{« $n \geqslant 2$ » en interligne au-dessus du $n \in \mathbb{N}^{*}$ biffé}
avec les relations 1.2. à 1.6. (où \struck{\ill{}} $x, y \in I$,
$\lambda \in A$, et $\pi^{1}(x)$ doit être lu $x$). \uncertain{Plus} \uncertain{aisé} :
façon de le dire est celle-ci : $I$ est un $A$-module,
on construit $\Gamma^{+}_{A}(\uncertain{I})$, et on
\note{la lettre dans $\Gamma^{+}_{A}(\ )$ est surchargée, $I$ ou $J$}

\page{24}
\uncertain{divise} ce \supplied{pseudo-}anneau à puissances divisées par l'idéal
engendré par les
\note{« pseudo- » ajouté en interligne}
\[
  (x.y)^{(n)} - x^{n} * y^{(n)}
\]
où le signe $*$ désigne \uncertain{provisoirement} la multiplication de
$\Gamma^{+}_{A}(\uncertain{I})$, le $.$ celle de $\uncertain{J}$. On constate que
les opérations $\pi^{n}$ passent au quotient $J'$ \note{« $J'$ » en interligne}
donnant une structure de $A$-algèbre à puissances divisées
\note{« de $A$-algèbre » en interligne} sur $J'$ [construction indépendante du
fait que $J$ est un idéal de $A$]. On
considère maintenant l'extension de $A$ par la $A$-algèbre $J'$
\note{« la $A$-algèbre » en interligne} définie à
l'aide de $J \to J'$, i.e. le quotient
\[
  A \times J' / \alpha(\uncertain{J}) , \quad \text{où} \quad \alpha : \uncertain{J} \to A \times J'
\]
est donné par
\[
  \alpha(x) = (x, -x')
\]
où $x \mapsto x'$ est l'application canonique $J \to J'$. \struck{On
vérifie \ill{} $(a, \xi) \in A \times J'$ \ldots\ $(a, \xi)(x, -x') = (ax, -\xi x')$}
\note{deux lignes biffées d'un trait, dont la fin — « $= ax, -\xi x'$ » — se lit ; un mot les suit, peut-être « \uncertain{sauf} »}

\emph{N.B.} Supposons que $A$ soit une algèbre sur un anneau $\Lambda$. Alors
l'hom composé $\Lambda \to A \to \Gamma(A, J)$ définit
évidemment sur $\Gamma(A, J)$ une structure de $\Lambda$-algèbre.
\page{25}
\emph{Cas particulier} \ : $A = \operatorname{Sym}_{\Lambda}(M)$ où $M$ est un
$\Lambda$-module, \struck{Alors} $J = \operatorname{Sym}^{+}_{\Lambda}(M)$ l'idéal
d'augmentation. Alors le pb universel est le même que celui
d'envoyer $M$ dans une $\Lambda$-algèbre à puissances divisées, on trouve
\[
  \Gamma\big(\operatorname{Sym}_{\Lambda}(M),\ \operatorname{Sym}^{+}_{\Lambda}(M)\big)
  \simeq \Gamma^{*}_{\Lambda}(M)
\]
Ainsi, on \uncertain{donne} \ill{} algèbre \struck{\ill{}} à puissances
divisées $\Gamma_{\Lambda}(M)$ définie comme foncteur
\ill{} les anneaux \emph{augmentés}
$\operatorname{Sym}^{*}_{\Lambda}(M)$ \ldots
\note{« $\Gamma^{*}_{\Lambda}(M)$ » : l'étoile est surchargée sur un $+$ ; deux mots sont biffés dans la ligne qui suit}

\struck{Si $A = A[T_{1}, \ldots, T_{n}]$}

[Soit $A$ un \struck{\ill{}} anneau avec idéal $J$ à puissances divisées,
on considère des éléments $(x_{i})_{i \in I}$ de $J$ et
proposons-nous de trouver les hom
$(A, J) \to (A', J')$ avec $A'$ anneau à idéal $J'$ à puissances
divisées, \note{« à puiss. div. » et « $J'$ » en interligne, « $A$-alg » biffé}
qui annulent les $x_{i}$. Foncteur \ill{} $(A', J')$ représentable ? Oui, \ill{}
la question de \uncertain{$A$} par \ill{}, il suffit de prendre
les $\pi^{n}(x_{i})$, $i \in I$, $n \geqslant 1$, et
l'idéal $K$ engendré par ces éléments, i.e. les
expressions $\sum \struck{\ill{}} \lambda_{n,i}\, \pi^{n}(x_{i})$ avec
$\lambda_{n,i} \in A$. On trouve ainsi le plus petit idéal $K$ stable par les
$\pi^{n}$ ($n \geqslant 1$),
\marginal{plus bas \ill{} si, aux conditions \ill{}}
\note{le début du crochet est un premier jet raturé : « $A$ » surchargé, « algèbre » biffé et remplacé en interligne par « anneau » ; dans la somme, un premier indice est biffé}

\page{26}
contenant les $x_{i}$, et pour un tel idéal $K$
(et pour un tel \uncertain{sous-anneau}), les $\pi^{n}$ passent au
quotient et définissent sur $(A/K, J/K)$ une structure d'anneau à idéal avec
puiss. divisées].
\note{« (et pour un tel \ldots) » en interligne}

Anneau avec idéal à puissances divisées \emph{nilpotentes} : qui, satisfait
\struck{$x^{(n)}$} $x^{(i)} = 0$ pour $i \geqslant \uncertain{n}$, $x \in J$,
ce qui implique pour $x \in J \Longrightarrow x^{\uncertain{n}} = 0$,
pour un ensemble fini d'éléments
\uncertain{ainsi} des nilpotents, \ill{} $J^{N+1} = 0$ ].
\struck{Application} \ill{} à puissances divisées « \struck{nilpotentes} »
\note{un pâté d'encre recouvre le $n$ dans « $i \geqslant n$ » ; l'exposant de $x^{n}$ est surchargé}

Si $J$ est engendré par $x_{1}, \ldots, x_{N}$, alors \ill{}
\emph{nilidéal} [de façon précise \ill{}
Anneau « avec » \emph{nilidéal} : \struck{$\forall x$, $\exists n$}
$\exists n$, tel que \struck{\ill{}} \ill{} à puissances
divisées nilpotentes. \struck{\ill{}} qu'il soit
\struck{\ill{}} à puissances divisées « nilpotentes »,
$\forall x$, $\exists n$ tel que $x^{(i)} = 0$ si $i > n$.
\note{trois formulations superposées et biffées en partie ; les deux définitions — nilidéal, puissances divisées nilpotentes — sont celles que le texte retient}

On a, pour un tel anneau, des
isomorphismes inverses l'un de l'autre
\[
  1 + J \;\rightleftarrows\; J, \qquad
  \log(1+x) = \sum_{n \geqslant 0} (-1)^{n}\,(n-1)!\; x^{(n)}, \qquad
  \exp x = \sum_{n \geqslant 0} x^{(n)}
\]
\note{la somme de $\log$ est écrite à partir de $n \geqslant 0$ ainsi sur la page, avec $(n-1)!$ ; la borne devrait être $n \geqslant 1$, et le signe attendu est $(-1)^{n-1}$}
\[
  1 + J \subset A^{*}
\]
\section{Site de De Rham}

\page{27}
\textbf{§ 4.} \emph{Site de De Rham}

\begin{tikzcd}[column sep=small, row sep=small]
U' \arrow[r] & U \arrow[r, hook] & X \arrow[d, "f"] \\
& & S
\end{tikzcd}
\note{le diagramme est tracé dans la marge gauche ; « $U \subset X$ » et la
flèche $U' \to U$ sont sur une même ligne, $f$ descend de $X$ vers $S$}

Schémas relatifs. Les objets du \emph{site de De Rham}
$(X/S)_{\mathrm{DR}}$ \struck{ou $\mathrm{DR}(X/S)$} sont les
systèmes \struck{triples} $(U, U', i, \pi)$ \struck{$(U, \mathcal{A}, \varphi, \pi^{*})$},
où $U$ est un ouvert de $X$, \struck{$\mathcal{A}$ un faisceau \ill{}
$i : U \to U'$} $U'$ un $S$-préschéma, $i : U \to U'$ une immersion
nilpotente, et $\pi$ une structure de puissances divisées \uncertain{topologiquement}
nilpotentes (\uncertain{pour} ou \uncertain{quasi-nilpotentes})
sur $\mathcal{J} = \operatorname{Ker}(\mathcal{O}_{U'} \to \mathcal{O}_{U})$.
Morphismes des sites \uncertain{évidents}.
\note{plusieurs mots-cadre des systèmes sont biffés et réécrits : un premier
quadruplet $(U, \mathcal{A}, \varphi, \pi^{*})$ et la phrase qui le glosait
sont barrés d'un trait, le quadruplet $(U, U', i, \pi)$ est écrit au-dessus ;
« topologiquement » est une lecture par le sens, le mot est abrégé}

Topologie \uncertain{ita} : la « topologie de Zariski ». Le site de De Rham
\uncertain{(strictifié)} est formé des objets \emph{au-dessus}
desquels \ill{} \uncertain{localement}
\ill{} \uncertain{$U'$} il existe une \uncertain{rétraction} de \uncertain{$U'$}
\ill{} (N.B. cette rétraction ne
fait pas partie de la structure).
\note{les deux lignes qui définissent le site strictifié sont d'une écriture
rapide ; « site de De Rham » et « (strictifié) » sont en interligne au-dessus
d'un mot biffé}

Topos $(X/S)^{\sim}_{\mathrm{DR}}$ ; $(X/S)^{\sim}_{\mathrm{DR\,st}}$ ou
\struck{\ill{}} $\widetilde{\mathrm{DR}}(X/S)$, $\widetilde{\mathrm{DR}}_{\mathrm{st}}(X/S)$ ;
si une ambiguïté n'est pas à craindre, on
laisse tomber le $/S$.

Diagramme de morphismes de topos

\page{28}
\begin{tikzcd}[column sep=small, row sep=small]
\widetilde{\mathrm{Zar}}(X) \arrow[r, dashed] & \widetilde{\mathrm{DR}}_{\mathrm{st}}(X) \arrow[r] \arrow[d] & \widetilde{\mathrm{Strat}}(X) \arrow[d] \\
& \widetilde{\mathrm{DR}}(X) \arrow[r] & \widetilde{\mathrm{Cris}}(X)
\end{tikzcd}
\note{un premier tracé de la ligne du haut, biffé et gribouillé, précède le
diagramme ; « st » de $\mathrm{DR}_{\mathrm{st}}$ est surchargé sur « strat ».
La flèche $\widetilde{\mathrm{Zar}}(X) \to \widetilde{\mathrm{DR}}_{\mathrm{st}}(X)$
est en pointillé sur la page}

Si $X$ est \uncertain{formellement} lisse sur $S$, alors les flèches
verticales sont des équivalences de Topos.

Si \struck{$X$} on est de car. 0, (p. ex.
$S$ de car. 0) alors les flèches horizontales sont des équivalences de Topos.

Dans les sites $\mathrm{DR}(X)$ et $\mathrm{DR}_{\mathrm{st}}(X)$, les produits
fibrés (de deux objets sur un troisième) et les noyaux de couples existent.
Mais il semble douteux que ce soit le
cas pour \uncertain{ceux} $\mathrm{Cris}(X)$ et $\mathrm{Strat}(X)$ \ldots ?
\note{le point d'interrogation est dans la marge, en regard de la dernière ligne}

Travaillons dans $\mathrm{DR}_{\mathrm{st}}(X)$. Notons que l'objet $\widetilde{X}$
de $\widetilde{\mathrm{DR}}_{\mathrm{st}}$ associé à $X$ comme
l'objet final (\ill{} \uncertain{définition} des
objets de $\mathrm{DR\,strat}$). Donc on peut utiliser
\struck{$\mathrm{DR\,strat}$} $\widetilde{X}$ pour écrire une
suite spectrale de Leray :
\[
  H^{*}(X_{\mathrm{DR\,st}},\ \mathcal{O}) \Longleftarrow
  E_{2}^{p,q} = H^{p}\big(\nu \mapsto H^{q}(\widetilde{X}^{\nu+1}, \mathcal{O})\big)
\]
\note{le nom de l'objet final est écrit $\widetilde{X}$ avec tilde ; l'indice
$\nu$ de la suite spectrale est une lecture, la lettre est petite}
\page{29}
Or $\widetilde{X}^{\nu+1}$ est \uncertain{ind-représentable} de la façon
suivante. On regarde les $X^{(\nu)} = X^{\nu+1}$, les voisinages
infinitésimaux $\Delta^{(\nu)}_{X/S}(i)$ de la diagonale. Pour chacun, on
regarde (qui à l'augmentation \uncertain{$\mathcal{O}_{\Delta^{(\nu)}(0)} = \mathcal{O}_{X}$})
le foncteur augmenté \uncertain{$\mathcal{O}_{\Delta^{(\nu)}(i)} \to \mathcal{O}_{\Delta^{(\nu)}(0)}$}
qu'il définit, sont à puissances divisées
\struck{\ill{}} c'est une notation d'ailleurs \struck{déjà} l'usage \ill{}
$\Delta^{(\nu)}_{X/S}(i)$ [\struck{\ill{}} justifiée, i.e. \uncertain{admettre}
\ldots\ déjà utilisé \ldots], et on y introduit
[pour $\mathcal{J}$ le \ill{} au \ill{} divisées] $\Delta^{(\nu)}_{X/S}(i)(j)$,
de sorte que
\[
  \Delta^{(\nu)}_{X/S}(i) = \varinjlim_{j} \Delta^{(\nu)}_{X/S}(i, j)
\]
\note{la ligne est écrite en deux fois, avec un premier
$\Delta^{(\nu)}_{X/S}(i)$ raturé ; les crochets sont de lui, et une partie
des mots qu'ils contiennent est biffée ou illisible}

\ill{} \uncertain{vue}
$\Delta^{(\infty)}_{X/S}(i)$ \note{souligné d'un trait ondulé} faisceau
\struck{De Rham} $\mathrm{DR}$-stratifiant « \struck{\ill{}}\uncertain{puissances} »
pour $\Delta^{(\nu)}_{X/S}(i)$ \ill{} \uncertain{ainsi}
$\Delta^{(\infty)}_{X/S}(i, \infty)$ ; \ill{}
\[
  \widetilde{X}^{\nu+1} = \varinjlim_{i, j} \Delta^{(\nu)}_{X/S}(i, j).
\]
\note{ce bloc est d'une écriture très rapide, avec plusieurs mots
surchargés ; les deux formules sont sûres, la prose entre elles ne l'est pas}

\marginal{à vérifier}
On aura encore
\[
  H^{q}\big(\mathrm{DR}_{\mathrm{st}}/(U, U') ;\ \mathcal{O}\big) \simeq H^{q}(U', \mathcal{O}_{U'})
\]
\note{« $\mathrm{DR}_{\mathrm{st}}$ » est surchargé sur un premier « DR » ;
le « à vérifier » est écrit en marge, en regard de la formule}

\page{30}
pour tout objet $U'$ du site $\mathrm{DR}$-stratifiant. D'où, si $X \to$
\uncertain{Aff}\ldots
\note{« DR- » en interligne au-dessus de « stratifiant » ; le mot après
$X \to$ est biffé}
\[
  \begin{cases}
    H^{q}(\widetilde{X}^{\nu+1}, \mathcal{O}) = 0 & \text{si } q > 0 \\
    H^{0}(\widetilde{X}^{\nu+1}, \mathcal{O}) = \varprojlim_{i,j}
      H^{0}\big(\Delta^{(\nu)}_{X/S}(i, j),\ \mathcal{O}_{\Delta^{\nu}_{X/S}(i,j)}\big)
  \end{cases}
\]
et par suite
\[
  H^{n}(X_{\mathrm{DR\,strat}}, \mathcal{O}) \simeq
  H^{n}\big(\nu \mapsto H^{0}(\widetilde{X}^{\nu+1}_{/X})\big)
\]
\note{le $H^{0}$ de droite est surchargé d'un $H^{p}$ ; l'indice $/X$ est
ainsi sur la page}

Dans le cas général, $X$ un anneau avec \ill{} puissances
divisées, \uncertain{sans} \ill{} \uncertain{considérer} $\mathcal{O}_{X}$ \ill{}
une \ill{} $\mathcal{E}^{*}$, et
\[
  H^{*}(X_{\mathrm{DR\,strat}}, \mathcal{O}) \simeq H^{*}(X_{\mathrm{Zar}}, \mathcal{E}^{*})
\]
\note{les deux lignes de prose qui introduisent $\mathcal{E}^{*}$ sont
d'une encre très pâle, presque effacée ; seul l'isomorphisme est sûr}
\section{L'algèbre universelle à puissances divisées sur un idéal}

\page{31}
Soit $A$ un anneau, $I$ idéal de $A$, \struck{$(x_{i})_{i \in I}$ un syst. de
gén. de $I$}. On cherche \struck{\ill{}} un $B$ avec idéal $J$ à puiss.
divisées, et hom $u : A \to B$ avec $u(I) \subset J$, qui soit
\emph{universel}. On prend l'algèbre quotient
\[
  B = A\Big[\struck{\ill{}}\ \{\Gamma^{n}(x)\}_{\substack{x \in I \\ n \geqslant 2}}\Big] \Big/ \text{relations}
\]
\note{le premier système de générateurs, entre le crochet ouvrant et l'accolade, est gribouillé jusqu'à l'illisible}
\begin{align*}
  \Gamma^{n}(x+y) &= \Gamma^{n}x + \Gamma^{n}y + \sum_{\substack{p, q \geqslant 1 \\ p+q = n}} \Gamma^{p}(x)\,\Gamma^{q}(y) \qquad (x, y \in I)
    \quad [\Gamma^{1}x = x \text{ par déf.}] \\
  \Gamma^{n}(xy) &= x^{n}\,\Gamma^{n}(y) \qquad x \in I,\ y \in A \\
  \struck{\Gamma^{m}(\Gamma^{n}(x))} &\struck{= \tfrac{(mn)!}{(n!)^{m}\, m!}\, \Gamma^{mn}(x)} \qquad x \in I \\
  \Gamma^{m}(x)\,\Gamma^{n}(x) &= \frac{(m+n)!}{m!\, n!}\; \Gamma^{m+n}(x)
\end{align*}
\note{les quatre relations sont accolées à gauche ; la troisième est gribouillée d'un trait épais et se lit sous les ratures ; « $\Gamma^{1}x = x$ par déf. » est encadré d'un trait à droite de la première}

$J$ idéal engendré par les $\Gamma^{n}(x)$ ($n \geqslant 1$, $x \in I$).
\struck{Supposons que \ill{}} À vérifier que c'est une solution du pb
universel. Il faut introduire sur $J$ des puiss. div. \ldots\ \struck{\ill{}}

\struck{$\Gamma_{A}$ \ill{} $\Gamma_{A}(\mathcal{J})$ \ill{} $A$-algèbre \ill{}
le quotient \ill{} $(xy)^{(n)} - x^{n} * y^{(n)}$ \ill{}
$*$ est \uncertain{provisoire}}
\note{un bloc de quatre lignes, encadré et hachuré de traits croisés ; il reprenait la construction de la page 24 — $\Gamma_{A}(J)$ divisé par les $(xy)^{(n)} - x^{n} * y^{(n)}$ — et n'en laisse lire que des bribes}

Supposons qu'on est dans un
cas où tout premier $\ell \neq p$
est inversible dans $A$, i.e. $A$ une
alg. sur $\mathbb{Z}_{p\mathbb{Z}}$ \note{« i.e. $A$ une alg. sur
$\mathbb{Z}_{p\mathbb{Z}}$ » en interligne, au-dessus d'un mot biffé}.
On peut simplifier la construction de $B$,
en utilisant que si $n$ est un entier, dont les
chiffres $p$-adiques \struck{\ill{}}
\[
  n = a_{0} + a_{1} p + a_{2} p^{2} + \cdots + a_{\nu} p^{\nu}, \qquad 0 \leqslant a_{i} \leqslant p-1,
\]
alors on aura
\[
  v_{p}(n) = \sum a_{i}\, v_{p}(p^{\nu}) \qquad \Big[\text{N.B. } v_{p}(p^{\nu}) = \frac{p^{\nu}-1}{p-1}\Big]
\]
\note{ainsi sur la page : $v_{p}(n)$ et $v_{p}(p^{\nu})$ sans factorielle, et l'indice $\nu$ sous la somme où l'on attendrait $i$ ; la formule que le N.B. donne est celle de $v_{p}(p^{\nu}\,!)$, et la ligne signifie $v_{p}(n!) = \sum a_{i}\, v_{p}(p^{i}\,!)$}
d'où si $C$ est une $A$-alg.
avec idéal $K$ à puiss. divisées, pour $x \in K$
\[
  \gamma^{n}(x) = c_{n}\; \gamma^{1}(x)^{a_{0}}\, \gamma^{p}(x)^{a_{1}} \cdots \big(\gamma^{p^{\nu}}(x)\big)^{a_{\nu}}
\]
\note{le second membre est écrit par-dessus un premier quotient gribouillé ; « $\gamma^{1}(x)^{a_{0}}$ » est surchargé}
où $c_{n} \in \mathbb{Z}_{p\mathbb{Z}}$ est donné par

\page{32}
\[
  c_{n} = \frac{\prod_{i} \big[(p^{i}\,!)/p^{\,v_{p}(p^{i}!)}\big]^{a_{i}}}{\big[n!/p^{\,v_{p}(n!)}\big]}
  = \text{partie première à } p \text{ de } \frac{\prod (p^{i}\,!)^{a_{i}}}{n!}
\]

Par suite, $B$ est aussi engendré par les éléments $\Gamma^{p^{\nu}}(x)$,
$\nu \geqslant 1$, $x \in I$, (quelles relations ?) qu'on
va noter $[x, \nu]$.

\emph{Définissons} les $\Gamma^{n}(x)$ dans $A\big[\{[x, \nu]\}_{\substack{x \in I \\ \nu \geqslant 1}}\big]$
(donc $\Gamma^{p^{\nu}}(x) = [x, \nu]$) par les
formules ci-dessus, en donnant les
relations \struck{\ill{}} \struck{$\Gamma^{p^{\nu}}$}
\note{le signe $\{$ devant le générateur est surchargé ; « relations » est écrit
au-dessus d'un mot biffé ; un $\Gamma^{p^{\nu}}$ gribouillé précède la
première formule}
\begin{align*}
  \Gamma^{p^{\nu}}(x+y) &= \Gamma^{p^{\nu}}x + \Gamma^{p^{\nu}}y
    + \sum_{\substack{\alpha, \beta \geqslant 1 \\ \alpha + \beta = p^{\nu}}} \struck{\ill{}}\ \Gamma^{\alpha}(x)\,\Gamma^{\beta}(y)
  \qquad \underbrace{\phantom{xxxxxxxxxxxxxx}}_{\text{expression en les } \Gamma^{p^{i}}(x),\ \Gamma^{p^{j}}(y),\ i, j < \nu} \\
  \Gamma^{p^{\nu}}(xy) &= x^{p^{\nu}}\, \Gamma^{p^{\nu}}(y) \\
  \Gamma^{p^{m}}(x)\,\Gamma^{p^{n}}(x) &= \frac{(p^{m} + p^{n})!}{p^{m}!\; p^{n}!}\; \Gamma^{p^{m}+p^{n}}(x) \qquad \ldots
\end{align*}
\note{les trois relations sont accolées à gauche ; dans la première un
coefficient est gribouillé avant $\Gamma^{\alpha}(x)\Gamma^{\beta}(y)$, et
l'accolade horizontale sous la somme porte la légende « expression en les
$\Gamma^{p^{i}}(x)$, $\Gamma^{p^{j}}(y)$, $i, j < \nu$ » ; dans la troisième
les exposants $m$ et $n$ sont lourdement surchargés et la formule se termine
par des points de suspension}

\marginal{ne \uncertain{faut-il} pas \uncertain{prendre}
les $\Gamma^{n}(x+y)$, toutes ?}

\marginal{il suffit de ? \ill{} en fait pour
\uncertain{certains} \ill{}. On a donc, si l'on
veut :
$\Gamma^{p^{\nu}}(x)^{p} = \dfrac{(p^{\nu+1})!}{(p^{\nu}!)^{p}\; p!}\; \Gamma^{p^{\nu+1}}(x)$}
\note{cette note marginale est écrite en oblique dans le coin inférieur gauche,
séparée du texte par un trait courbe ; la formule est le cas $m = p$, $n = p^{\nu}$
de 1.4}

Supposons maintenant que $I$ soit engendré par un élément $\lambda$. Alors $B$
est aussi isomorphe au quotient de
$A\big[(\Gamma^{n}(\lambda))_{n \geqslant 2}\big]$ par les relations
\note{la phrase s'arrête là, en bas de la page ; sa suite n'est pas dans ce
lot}
\page{34}
\note{le haut de la page continue la phrase interrompue au bas de la page 32 : le quotient de $A[(\Gamma^{n}(\lambda))_{n \geqslant 2}]$ par les relations qui suivent}
\[
  \begin{cases}
    \Gamma^{m}(\lambda)\,\Gamma^{n}(\lambda) - \dfrac{(m+n)!}{m!\, n!}\; \Gamma^{m+n}(\lambda) \\[6pt]
    \struck{\ill{}} \quad (x^{n} - y^{n})\,\Gamma^{n}(\lambda) \qquad \text{pour } x, y \in A \text{ tels que } (x-y)\lambda = 0 .
  \end{cases}
\]
\note{la seconde ligne de l'accolade commence par une expression gribouillée jusqu'à l'illisible ; « tels que » est en interligne}

[\emph{N.B.} Ces relations \uncertain{disent} que
$\Gamma^{n}(\lambda x)$ doit être $\equiv x^{n}\,\Gamma^{n}(\lambda)$,
\ill{} \struck{\ill{}} \ill{} \uncertain{en} $\lambda x$
i.e. \struck{\ill{}} si $\lambda x = \lambda y$
alors $x^{n}\,\Gamma^{n}(\lambda) \equiv y^{n}\,\Gamma^{n}(\lambda)$
\marginal{qui : $(**)$}, les relations
$\Gamma^{n}(\lambda(x+y)) \equiv \cdots$ résulte alors de
$(**)$, et $\Gamma^{n}(xy) \equiv x^{n}\,\Gamma^{n}(y)$ est
\uncertain{trivial}, ainsi que $\Gamma^{m}(\lambda x)\,\Gamma^{n}(\lambda y) = \cdots$ ]
\note{le N.B. est encadré de crochets sur la page ; « qui : $(**)$ » est écrit en marge à droite, cerclé, en regard de la relation $x^{n}\Gamma^{n}(\lambda) \equiv y^{n}\Gamma^{n}(\lambda)$, et c'est le nom qu'il lui donne}

Si maintenant tous les nombres premiers $\neq p$ sont
inversibles dans $A$, on trouve aussi le quotient de
\[
  A\big[(\Gamma^{p^{\nu}}(\lambda))_{\nu \geqslant 1}\big]
\]
par les relations
\[
  \begin{cases}
    \Gamma^{m}(\lambda)\,\Gamma^{n}(\lambda) - \dfrac{(m+n)!}{m!\, n!}\; \Gamma^{m+n}(\lambda) \\[6pt]
    (x^{n} - y^{n})\,\Gamma^{n}(\lambda) \qquad \text{si } x, y \in A \text{ tels que } (x-y)\lambda = 0
  \end{cases}
\]
où $\Gamma^{n}(\lambda)$ se définit en termes
des $\Gamma^{p^{\nu}}(\lambda)$ comme plus haut. Les
premières relations se réduisent à
\[
  \boxed{\ \Gamma^{p^{\nu}}(\lambda)^{\uncertain{e}} = \struck{\ill{}}\; \frac{(\uncertain{e}\,p^{\nu+1})!}{(p^{\nu}!)^{p}}\; \Gamma^{\uncertain{e}\,p^{\nu+1}}(\lambda)\ }
\]
pour $p \leqslant e < 2p-1$, donc pour $e = p + a$ \ ($0 \leqslant a \leqslant p-1$).
\note{la formule encadrée est lourdement surchargée : l'exposant de $\Gamma^{p^{\nu}}(\lambda)$, d'abord $p$ (comme dans la note marginale de la page 32), est réécrit en un signe plus épais, et une première fraction $(a+b)!/\ldots$ est gribouillée devant celle qui reste. Telle qu'elle se lit, la formule ne suit pas de 1.5 : pour l'exposant $e$ on attendrait $(e p^{\nu})!/(p^{\nu}!)^{e}\; \Gamma^{e p^{\nu}}(\lambda)$. Elle est reproduite telle quelle}

\page{36}
\emph{Lemme de Poincaré} vrai pour les $k$-algèbres suivantes
($k$ corps de car. nulle) :
\[
  \begin{cases}
    k[[x]][Y] / (Y^{p} - x^{q}) \\
    k[[x]][Y] / (Y^{2} + P(x)\,Y + Q(x)) \\
    k[x_{1}, \ldots, x_{n}] / (x_{1}^{p_{1}}, x_{2}^{p_{2}}, \ldots, x_{n}^{p_{n}}) \qquad \text{p. ex. } k[x]/(x^{p}) \\
    k[x, y] / \mathfrak{m}^{2} \\
    k[x, y] / (x^{3}, y^{2}, x^{2}y) \\
    k[x, y] / (x^{3}, y^{2}, xy)
  \end{cases}
\]
\note{le $Y$ de la première ligne est réécrit par-dessus un $X$ ; dans la
troisième, les exposants $p_{i}$ sont ainsi sur la page, et ne sont pas le $p$ des
pages précédentes}

Si $A$ \struck{\ill{}} \uncertain{local} sur $k$, ext. résiduelle triviale, alors
pour tout $x \in \mathfrak{m}$ avec $dx = 0$ dans $\Omega$,
\ldots\ $x \in \mathfrak{m}^{3}$.
\note{le mot biffé après $A$ est réécrit en interligne « local » ; « dans $\Omega$ » est sous le $dx = 0$ ; le membre de droite se lit $x \in \mathfrak{m}^{3}$, l'exposant est sûr}

\marginal{Ann. Éc. N. Sc. 65 ou 66, J. de Liouville}
\note{la référence est de sa main, sans autre précision ; le « 65 ou 66 » est encadré}

\page{38}
\emph{Supposons que les cas. \uncertain{résiduelles} de $A$ soient \uncertain{parfaites} : $p > 0$}
\struck{$(x^{(m)})^{(n)} =$}\quad opérations de puiss. div. $(\pi_{m})_{m \geqslant 0}$
\note{l'intitulé est souligné ; « $(x^{(m)})^{(n)} =$ » est gribouillé sous lui ; le
« cas. » abrège sans doute « caractéristiques » ; les opérations sont notées
ici $\pi_{m}$ en indice, et non plus $\pi^{m}$ en exposant comme aux pages 21–26}

\begin{enumerate}
  \item[a)] $\pi_{m}(x)\,\pi_{n}(x) = \binom{m+n}{n}\, \pi_{m+n}(x)$
  \note{les indices de la formule a) sont lourdement surchargés : un premier
  $\pi^{(m)}(x)\,\pi^{(n)}(x) = \ldots\, \pi^{(m+n)}(x)$ en exposants est
  réécrit avec les indices en bas}

  Supposons que $\begin{cases} m \equiv 0 \ (p) \\ 0 \leqslant n \leqslant p-1 \end{cases}$
  alors \struck{\ill{}} $\binom{m+n}{n} = \dfrac{(m+1)\cdots(m+n)}{n!}$
  \struck{\ill{}} est inversible \uncertain{(p)}, donc
  \[
    \boxed{\ \pi_{m+n}(x) = \binom{m+n}{n}^{-1}\, \pi_{m}(x)\,\pi_{n}(x)\ }
  \]
  \note{la formule encadrée est surchargée elle aussi : les indices $m+n$, $m$, $n$
  sont réécrits par-dessus des exposants entre parenthèses}

  \item[b)] $\pi_{m} \circ \pi_{n} = \dfrac{(mn)!}{(n!)^{m}\, m!}\; \pi_{mn}$

  [Supposons $m = p$, on trouve un coefficient
  \[
    c_{n} = \frac{(pn)!}{(n!)^{p}\; p!}
  \]
  et on a, si $n = a_{0} + a_{1} p + \cdots + a_{\nu} p^{\nu}$ :
  \begin{align*}
    v_{p}(c_{n}) &= v_{p}\big((pn)!\big) - p\, v_{p}(n!) - 1 \\
      &= \frac{pn - \sum a_{i}}{p-1} - p\,\frac{n - \sum a_{i}}{p-1} - 1 \\
      &= \struck{\ill{}}\ \Big(\sum a_{i}\Big) - 1
  \end{align*}
  \note{dans la dernière ligne un premier coefficient devant la somme est
  gribouillé ; le calcul utilise $v_{p}(n!) = (n - \sum a_{i})/(p-1)$}

  Supposons $n = p$, on trouve un coeff.
  \[
    d_{m} = \frac{(pm)!}{(p!)^{m}\; m!}
  \]
  et, si $m = b_{0} + b_{1} p + \cdots + b_{s} p^{s}$,
  \begin{align*}
    v_{p}(d_{m}) &= v_{p}\big((pm)!\big) - m\, v_{p}(p!) - v_{p}(m!) \\
      &= \frac{pm - \sum b_{i}}{p-1} - m - \frac{m - \sum b_{i}}{p-1} \\
      &= 0
  \end{align*}
  donc
  \[
    \pi_{pm} = \Big(\frac{(pm)!}{(p!)^{m}\, m!}\Big)^{-1}\, \pi_{m} \circ \pi_{p}\ ]
  \]
  \note{le crochet ouvert avant « Supposons $m = p$ » se ferme après cette
  formule ; devant $m = b_{0} + \cdots$ un $\sum$ est biffé}
\end{enumerate}

a) et b) ensemble montrent que les $\pi_{n}$
sont connus \emph{quand on connaît} $\pi_{p}$ \note{ainsi souligné}

\page{40}
$\pi = \pi_{p}$ satisfait les conditions
\[
  \begin{cases}
    \pi(x+y) = \pi x + \pi y + \displaystyle\sum_{0 < i < p} \struck{\ill{}}\ \frac{1}{i!\,(p-i)!}\; x^{i} y^{p-i} \\[8pt]
    \pi(fx) = f^{p}\, \pi x
  \end{cases}
\]
\note{un premier coefficient devant la somme est gribouillé ; l'exposant de
$\pi$ dans $\pi(x+y)$ est surchargé ; le $y^{p-i}$ se lit $y^{i}$ avec un
exposant corrigé}

Si on a une opération
satisfaisant ces conditions, peut-on
en définir les $\pi^{(n)}$ \struck{\ill{}}
satisfaisant les conditions $(1.1)$ à $(1.5)$
\struck{\ill{}}, avec $\pi^{(p)} = \pi$ \struck{Par récurrence} des
puiss. div. \note{les mots entre $\pi^{(n)}$ et
« satisfaisant » sont surchargés ; « des puiss. div. » est écrit sous « Par
récurrence » biffé}

\struck{On va construire les $\pi_{N}$ par récurrence, en
supposant les $\pi_{i}$ construits pour $i < N$ et \uncertain{façon}
à satisfaire $(1.1)$ à $(1.5)$ [ou $1.\uncertain{4}$
pour $mn < N$, dans $1.5$ pour $m+n < N$]. Cas $1$ : $N \not\equiv 0\ (p)$,
donc $N = N' + s$, avec $N' \equiv 0\ (p)$, $0 < s \leqslant p-1$. On aura donc \ldots}
\note{tout ce paragraphe est barré de deux longs traits obliques et d'un trait
courbe qui le contourne ; le « Cas 1 » et la décomposition $N = N' + s$ se lisent
sous les traits ; « $0 < s \leqslant p-1$ » est en interligne sous « avec »}

On va construire
\[
  \pi_{n}(x) = c_{n}\; x^{a_{0}}\, \pi(x)^{a_{1}}\, \pi^{2}(x)^{a_{2}} \cdots \pi^{r}(x)^{a_{r}}
\]
où $n = a_{0} + a_{1} p + \cdots + a_{r} p^{r}$ est le
développement $p$-adique de $n$, et où $c_{n} \in \struck{\mathbb{Z}_{p}\mathbb{Z}^{*}}\;
\mathbb{Q} \cap \mathbb{Z}_{p}^{*}$ et un élément
de $\mathbb{Q}$ dont la valuation $p$-adique est nulle. [Il
serait intéressant de connaître ce
$c_{n}$ mod $p$, i.e. un élément de $\mathbb{F}_{p}^{*}$.
Est-ce tjs $1$ ?]
\note{$\pi^{2}(x)$, $\pi^{r}(x)$ désignent ici les itérés $\pi \circ \pi$, \ldots, de
$\pi = \pi_{p}$ ; le premier « $\mathbb{Z}_{p}\mathbb{Z}^{*}$ » est biffé et
remplacé par $\mathbb{Q} \cap \mathbb{Z}_{p}^{*}$ ; la question finale est de lui.
La suite de ce calcul n'est pas dans ce lot}

\end{document}
